%% ============================================================
%%  ML peak finder for VFE board QC voltage traces
%%  Informal lab talk, Beamer / pdflatex, dark-beamer style
%% ============================================================
\documentclass[aspectratio=169,10pt]{beamer}

\usepackage[T1]{fontenc}
\usepackage[utf8]{inputenc}
\usepackage{lmodern}
\usepackage{graphicx}
\usepackage{booktabs}
\usepackage{array}
\usepackage{xcolor}
\usepackage{amsmath}
\usepackage{amssymb}
\usepackage{microtype}
\usepackage{hyperref}
\usepackage{tikz}
\usetikzlibrary{arrows.meta, positioning, fit, backgrounds, calc}
\usetikzlibrary{decorations.pathmorphing, decorations.markings}
\usepackage{pgfplots}\pgfplotsset{compat=1.16}

\graphicspath{{../plots/}}   % <-- point at your figure directory

% ── Colour palette (CMS dark-navy style) ─────────────────────
\definecolor{navydark}  {RGB}{10,  27, 56}
\definecolor{navymid}   {RGB}{18,  48, 96}
\definecolor{cmsblue}   {RGB}{ 0, 112,192}
\definecolor{cmsaccent} {RGB}{ 0, 188,235}
\definecolor{cmsgold}   {RGB}{255,185,  0}
\definecolor{lightgrey} {RGB}{210,220,235}
\definecolor{offwhite}  {RGB}{240,244,250}

\newcommand{\shorttitletext}{ML peak finder for VFE QC traces}

% ── Beamer theme setup ───────────────────────────────────────
\usetheme{default}
\setbeamertemplate{navigation symbols}{}
\setbeamertemplate{itemize item}{\color{cmsaccent}\small$\blacktriangleright$}
\setbeamertemplate{itemize subitem}{\color{cmsgold}\footnotesize$\bullet$}
\setbeamercolor{background canvas}{bg=navydark}
\setbeamercolor{normal text}{fg=offwhite}
\setbeamercolor{alerted text}{fg=cmsgold}
\setbeamercolor{example text}{fg=cmsaccent}
\setbeamercolor{frametitle}{fg=cmsaccent, bg=navymid}
\setbeamerfont{frametitle}{size=\large, series=\bfseries}
\setbeamertemplate{frametitle}{%
  \vskip2pt
  \hspace*{6pt}\insertframetitle
  \vskip3pt
  \color{cmsaccent}\rule{\linewidth}{0.4pt}%
  \vskip2pt
}
\setbeamercolor{structure}{fg=cmsaccent}
\setbeamercolor{item}{fg=cmsaccent}
\setbeamercolor{subitem}{fg=cmsgold}
\setbeamercolor{footline}{fg=lightgrey, bg=navymid}
\setbeamertemplate{footline}{%
  \begin{beamercolorbox}[wd=\paperwidth,ht=2.2ex,dp=1ex]{footline}
    \hspace*{6pt}\tiny \shorttitletext
    \hfill \tiny\insertframenumber\,/\,\inserttotalframenumber\hspace*{6pt}
  \end{beamercolorbox}%
}
\setbeamercolor{block title}{fg=navydark, bg=cmsaccent}
\setbeamercolor{block body}{fg=navydark, bg=offwhite}
\setbeamercolor{block title alerted}{fg=navydark, bg=cmsgold}
\setbeamercolor{block body alerted}{fg=navydark, bg=offwhite}
\hypersetup{colorlinks=true, urlcolor=cmsaccent, linkcolor=cmsaccent}

% ── Macros ───────────────────────────────────────────────────
\newcommand{\pT}{\ensuremath{p_{\mathrm{T}}}}
\newcommand{\etal}{\ensuremath{|\eta|}}
\newcommand{\Hgg}{\ensuremath{H\!\to\!\gamma\gamma}}
\newcommand{\Zee}{\ensuremath{Z\!\to\!ee}}
\newcommand{\RNine}{\ensuremath{R_{9}}}
\newcommand{\hl}[1]{\textcolor{cmsaccent}{\textbf{#1}}}
\newcommand{\gold}[1]{\textcolor{cmsgold}{\textbf{#1}}}
\newcommand{\plotslide}[3][0.80]{%
  \begin{center}
    \includegraphics[width=#1\linewidth,height=0.74\textheight,keepaspectratio]{#2}\\[2pt]
    {\tiny\color{lightgrey}#3}
  \end{center}}

% ── White frame (for CMS public plots on white canvas) ───────
\newenvironment{whiteframe}[1]{%
  \setbeamercolor{background canvas}{bg=white}%
  \setbeamercolor{normal text}{fg=navydark}%
  \setbeamercolor{frametitle}{fg=navymid, bg=white}%
  \setbeamertemplate{frametitle}{\vskip2pt\hspace*{6pt}\insertframetitle\vskip3pt
    \color{navymid}\rule{\linewidth}{0.4pt}\vskip2pt}%
  \setbeamercolor{footline}{fg=navymid, bg=white}%
  \usebeamercolor[fg]{normal text}%
  \begin{frame}{#1}}{\end{frame}}
\newcommand{\whitecap}[1]{{\tiny\color{navymid}#1}}
\newcommand{\TODO}[1]{\textcolor{cmsgold}{$\langle$\,TODO: #1\,$\rangle$}}
\newcommand{\partdivider}[2]{%
  {\setbeamercolor{background canvas}{bg=navymid}
   \setbeamertemplate{footline}{}
   \begin{frame}[plain]
     \vspace*{2.0cm}\centering
     \ifx&#1&\else{\large\color{lightgrey} #1}\\[8pt]\fi
     {\huge\bfseries\color{cmsaccent} #2}
   \end{frame}}}



\begin{document}

% ── Title page ───────────────────────────────────────────────
{
\setbeamertemplate{footline}{}
\begin{frame}[plain]
  \vspace*{1.2cm}
  \begin{center}
    {\color{cmsaccent}\rule{0.6\linewidth}{1.5pt}}\\[10pt]
    {\Large\bfseries\color{white} Finding P0--P4 in VFE quality-control traces}\\[6pt]
    {\Large\bfseries\color{white} with a learned peak finder}\\[10pt]
    {\color{cmsaccent}\rule{0.6\linewidth}{1.5pt}}\\[20pt]
    {\large\color{white} Santanu}\\[10pt]
    {\footnotesize\color{cmsgold} informal lab discussion --- 10 November 2026}\\[6pt]
    {\footnotesize\color{lightgrey} status: toy study only, no board data yet}
  \end{center}
\end{frame}
}

% ── 1. The problem ───────────────────────────────────────────
\begin{frame}{Where this comes from}
  \begin{itemize}
    \item VFE board QC gives a \hl{voltage trace} per board. Five landmarks have
          to be read off it: \gold{P0, P1, P2, P3, P4}.
    \item Today a \hl{hard-coded finder} does this --- fixed thresholds and
          windows.
    \item It \gold{misidentifies peaks} on some boards. Every miss needs a human
          to look at the trace.
    \item Idea: learn the peak positions from labelled traces instead of
          hand-tuning thresholds, and get a \hl{confidence per peak} so the
          doubtful boards flag themselves.
  \end{itemize}
  \vfill
  {\tiny\color{lightgrey} Everything shown here is measured on simulated traces built from the
   description of the real waveform --- see the next slide.}
\end{frame}

% ── 2. The waveform ──────────────────────────────────────────
\begin{frame}{What the trace looks like}
  \plotslide[0.92]{anatomy.png}{One simulated trace. Gold markers are the five
  true landmark positions; P0 and P1 sit close together, 10--60 samples apart.
  Amplitude is in arbitrary units, shape only --- the absolute scale is not
  modelled.}
\end{frame}

% ── 3. Why it is not a threshold problem ─────────────────────
\begin{frame}{Why thresholds struggle}
  \begin{itemize}
    \item \gold{P0 hides in the spike train.} It is the \emph{last and smallest}
          of many little spikes --- nothing about its height singles it out.
    \item \gold{P0 then sits right in front of P1}, a few tens of samples away
          and far smaller. A finder that keys on height takes the pair as one
          structure.
    \item \gold{The floor never resets.} After each peak the baseline settles
          \hl{higher} than before, so one global threshold cannot serve all five.
    \item \gold{There is structure that is not a peak.} Wiggles ride on the
          raised floors and must be ignored.
    \item \gold{The shape varies.} Rise and fall widths are drawn independently:
          a peak can rise slowly and fall fast, or the other way round.
    \item Positions are not fixed --- the five landmarks move trace to trace.
  \end{itemize}
\end{frame}

% ── 4. Method ────────────────────────────────────────────────
\begin{frame}{The approach}
  \vspace{0pt}\par
  {\small A fully convolutional network reads the whole trace and emits
   \hl{five output channels}, one per landmark.}
  \par\vspace{4pt}
  \begin{itemize}
    \item Each channel is trained to peak at its own landmark --- a narrow bump,
          width \hl{4 samples}, centred on the true position.
    \item Position $=$ where the channel is largest. \gold{Confidence $=$ how
          large it gets} --- that is the per-peak quality flag.
    \item Channels are separate, so \gold{P0 cannot be confused with P4}: each
          one only has to find its own landmark.
    \item Each landmark is reported with \gold{the apex beside it} --- the
          nearest local maximum and its height, in the board's own units. If
          that apex is clearly higher than the trace at the reported position
          and the channel is unsure, the position moves onto it.
    \item Trained on \hl{8000 simulated traces}, tested on 1000 held out.
          Input is median-subtracted and scaled per trace, so absolute gain
          does not matter.
  \end{itemize}
  \vfill
  {\tiny\color{lightgrey} Alternative considered: regress five numbers directly. Rejected ---
   it gives no per-peak confidence, which is the part QC actually needs.}
\end{frame}

% ── 5. It works on toys ──────────────────────────────────────
\begin{frame}{One test trace}
  \plotslide[0.92]{heatmap_demo.png}{Top: a held-out trace, true positions
  dashed, network output solid. Bottom: the five channel outputs, with the
  confidence each one reaches. Amplitude normalised per trace.}
\end{frame}

% ── 6. Numbers ───────────────────────────────────────────────
\begin{frame}{How well it does --- on simulation}
  \vspace{0pt}\par
  {\small All five within 2 samples on \gold{76.8\%} of traces, within 4 on
   \gold{89.0\%}. Badly wrong ($>15$ samples) on \gold{2.0\%}.}
  \par\vspace{2pt}
  \plotslide[0.88]{perpeak.png}{1000 held-out simulated traces. Fraction of
  traces where each landmark lands within 2 (left) and 4 (right) samples of the
  visible apex. Suppressed zero on both.}
  \par\vspace{-4pt}
  {\tiny\color{lightgrey} Labels sit on the visible apex. Misses are mostly 3 samples --- just
   over the line, no early/late bias. On P3 and P4 most of them land on a flat top within
   1\% of the apex height.}
\end{frame}

% ── 7. The hybrid, and why it is not the default ─────────────
\begin{frame}{A hybrid I tried --- and dropped}
  \vspace{0pt}\par
  {\small Two-stage idea: scan the trace for local maxima, then let a classifier
   decide which candidate is which landmark.}
  \par\vspace{2pt}
  \plotslide[0.86]{methods.png}{Left: fraction of traces with all five landmarks
  within 2 samples. Right: how often the true landmark is among the scanned
  candidates, against the scan threshold; the label gives candidates per trace.}
  \par\vspace{-4pt}
  {\tiny\color{lightgrey} With labels on the visible apex the scan finds the true peak 99.9\%
   of the time, so recall is not the issue --- the re-ranker is simply a worse judge than
   the network it is reading.}
\end{frame}

% ── 8. Next ──────────────────────────────────────────────────
\begin{frame}{What I need next}
  \begin{itemize}
    \item \gold{Real traces.} Everything above is simulation written from a
          verbal description --- the numbers are a sanity check, not a result.
    \item \gold{Labels on the hard cases} --- especially traces the present
          finder gets wrong. A few hundred is enough to fine-tune.
    \item \gold{One detail to check:} a peak is its \hl{visible apex} --- but
          does the current finder take the apex of the \hl{raw} trace or of a
          smoothed one? Against the raw-trace apex, snapping to it gains
          \gold{20 points}.
    \item Then: refit the simulation to measured traces, fine-tune, and report
          this same table on real held-out boards.
  \end{itemize}
\end{frame}

% ── Backup ───────────────────────────────────────────────────
\partdivider{}{Backup}

\begin{frame}{Backup: more simulated traces}
  \plotslide[0.95]{gallery.png}{Six independent draws, identical generator
  settings. Gold markers are the true landmarks. Each panel is auto-scaled;
  shape only, no common vertical scale.}
\end{frame}

\begin{frame}{Backup: how a trace is built}
  \vspace{0pt}\par
  {\footnotesize Every number below is a guess from the verbal description.
   \gold{These are what real traces will replace.}}
  \par\vspace{3pt}
  \centering\footnotesize
  \begin{tabular}{@{}l l@{}}
    \toprule
    \color{cmsaccent}\textbf{ingredient} & \color{cmsaccent}\textbf{range used} \\
    \midrule
    trace length                       & 1024 samples \\
    P1--P4 live in                     & samples 143--922, evenly spread \\
    P0 apex ahead of P1 apex by        & 10--60 samples \\
    spikes before P0                   & 8--20, height 0.01--0.06 \\
    P0 height                          & 0.04--0.09 (top of the spike family) \\
    P1--P4 heights                     & 0.25--1.00, increasing within a trace \\
    leading-edge width                 & 1.5--14 samples \\
    trailing decay                     & 2.5--30 samples, drawn independently \\
    per-trace shape tilt               & 0.5--2.0 (slow rise $\leftrightarrow$ fast rise) \\
    fraction of a peak left as floor   & 0.25--0.55 \\
    floor step width                   & 3--10 samples \\
    decoy bumps per floor              & 2--6, at 4--18\% of the local step \\
    noise                              & 0.004--0.015, plus 0.2\% single-sample glitches \\
    label                              & visible apex; 42\% of draws redrawn \\
    landmark spacing jitter            & 50\% of the mean spacing \\
    \bottomrule
  \end{tabular}
\end{frame}

\begin{frame}{Backup: network and training}
  \begin{itemize}
    \item \hl{U-shaped 1D convolutional network} --- three downsampling levels,
          three up, skip connections, kernel width 9. \gold{462k parameters}.
    \item Output: five channels at full trace length. Target is a Gaussian of
          width \hl{4 samples} on the true position; loss is per-sample binary
          cross-entropy against it.
    \item Adam, learning rate $10^{-3}$ with cosine decay, batch 32,
          \hl{25 epochs}. \gold{125 s} on one GPU.
    \item Input per trace: subtract the median, divide by the largest absolute
          value. \gold{Any inference path must repeat this exactly} --- it is
          part of the model.
    \item Labels and errors use the \hl{visible apex} of the noiseless
          waveform. Draws where a landmark has no visible apex are redrawn.
  \end{itemize}
\end{frame}

\begin{frame}{Backup: where the failures are}
  \plotslide[0.78]{p0_errors.png}{Held-out simulated traces: P0 (1000 entries,
  filled) against P1--P4 pooled (4000 entries, outline). Log vertical scale,
  dotted line at the 2-sample criterion, everything beyond 20 piled in the last
  bin.}
  \par\vspace{-4pt}
  {\tiny\color{lightgrey} P0 misses on 40 traces, 8 badly. About half land earlier, in the
   spike train; none land on P1.}
\end{frame}

\begin{frame}{Backup: snapping to the apex}
  \plotslide[0.88]{snap.png}{Left: moving each reported position onto the nearby
  apex of the noisy trace, against how many positions get moved. Right: the same
  two methods scored against the clean apex and against the noisy-trace apex.
  1000 held-out simulated traces.}
  \par\vspace{-4pt}
  {\tiny\color{lightgrey} Scored against the clean apex, snapping does nothing. Scored against
   the apex of the noisy trace, always snapping gains \hl{20 points} --- so it matters which
   of the two the present finder reports.}
\end{frame}

\begin{frame}{Backup: the hybrid in numbers}
  \vspace{0pt}\par
  {\footnotesize Same 1000 held-out traces; the re-ranker was trained on 2000
   separate ones.}
  \par\vspace{3pt}
  \centering\small
  \begin{tabular}{@{}l c c@{}}
    \toprule
    \color{cmsaccent}\textbf{method} & \color{cmsaccent}\textbf{all five within 2}
      & \color{cmsaccent}\textbf{badly wrong} \\
    \midrule
    network alone                  & \gold{76.8\%} & \gold{2.0\%} \\
    scan + ranker, ordered decode  & 71.2\%        & 2.1\% \\
    scan + ranker, plain decode    & 71.4\%        & 2.1\% \\
    scan + ranker, greedy          & 71.2\%        & 2.4\% \\
    \bottomrule
  \end{tabular}
  \par\vspace{6pt}
  \begin{itemize}\footnotesize
    \item Ordering constraints are \gold{not} the problem: ordered and plain
          decode agree to a fifth of a percent, and the network's own answer
          breaks the ordering on only \hl{1 of 1000} traces.
    \item After adding the network's own maxima to the candidate list, the true
          landmark is in the list on \hl{1000 of 1000} traces --- so recall is
          not the problem either.
    \item Conclusion: the second-stage ranker is a \gold{worse judge than the
          network output it reads}. Kept as a switch, not the default.
  \end{itemize}
\end{frame}

\end{document}
